```latex
\section{Introduction: From ``It's the Vibe'' to Vibe Research}

In the Australian film \emph{The Castle}, the inexperienced lawyer Dennis Denuto is convinced that his clients have a valid constitutional case, but struggles to express the underlying principle in formal legal terms. His argument eventually collapses into the memorable phrase, ``It's the vibe.'' Only when an experienced barrister reformulates the intuition in language recognised by the court does the case succeed.

The expression \emph{vibe coding} began with something of the same deliberately informal quality. Introduced by Andrej Karpathy in February 2025, it named a style of AI-mediated programming whose practical usefulness was becoming apparent before a settled conceptual vocabulary had developed around it \citep{fowler2026vibecoding}. What initially sounded like a joke has rapidly become a recognisable software-development methodology.

In a contemporary vibe-coding workflow, a human developer may formulate requirements, inspect visible behaviour, request modifications, perform selected independent checks, and ultimately decide whether a system is ready to ship. An advanced coding assistant may meanwhile perform much of the detailed implementation, debugging, refactoring, test construction, repository inspection, and technical explanation.

The point is not that such development is risk-free. It plainly is not. A developer who delegates substantial work to an AI system must decide how much independent checking is appropriate. With experience, this becomes partly a problem of calibrated trust: learning which outputs require close inspection, what kinds of tests provide useful evidence, which failures are characteristic of the system being used, and when additional verification is needed.

One of the principal attractions of contemporary vibe coding is precisely that sufficiently capable coding assistants make exhaustive line-by-line human inspection unnecessary in many cases. If every generated component had to be independently reconstructed or checked in detail, much of the productivity advantage would disappear.

Nor does the supervising human necessarily acquire detailed knowledge of everything the coding assistant has done. Expertise may be distributed between the human and the AI system. If another team member asks why an obscure component was implemented in a particular way, the most useful response may be to ask the coding assistant to inspect the repository and explain it.

These features suggest a natural question. What happens when essentially the same methodology is applied to academic research?

Advanced AI systems can now participate in formulating research questions, reviewing literature, proposing hypotheses, designing experiments, implementing analyses, interpreting results, constructing arguments, drafting papers, and responding to criticism. We use the term \emph{vibe research} for sustained iterative research collaboration of this kind.

Gunkel approaches human--AI collaboration from a very different direction. Drawing on Barthes, Derrida, and the history of distributed creative production, he challenges both the instrumental conception of generative AI as a passive tool and the romantic conception of the author as a singular origin of meaning; the essay is itself presented as the product of an iterative collaboration with ChatGPT-4o \citep{gunkel2025prompted}. Our starting point is instead software-engineering practice, in particular the delegation of substantive work to agents operating within versioned repositories. It is striking that these philosophical and technical routes lead to a related conclusion: the model of a sovereign human author employing a passive instrument describes the observed process poorly, and intellectual production is better understood as distributed, iterative, and dialogical.

The central puzzle of this paper follows directly. If vibe coding and vibe research involve closely similar forms of human--AI collaboration, why has vibe coding become an increasingly accepted development practice while vibe research remains difficult to accommodate within mainstream academic norms?

Arguably the most important difference is sociological. Research does not consist only in producing intellectual artefacts. Being a researcher normally means participating in a research community: publishing work under one's name, responding to questions and criticism, explaining how results were obtained, and accepting various forms of responsibility for what one has helped produce.

This is often phrased as some version of the question \emph{``Can the AI be a responsible author?''}, with the imputation that it cannot. In current mainstream academic practice, AI authors are generally not permitted. The International Committee of Medical Journal Editors, for example, states that chatbots should not be listed as authors because they cannot be responsible for the accuracy, integrity, and originality of the work \citep{icmje2026ai}. Nature Portfolio similarly states that authorship carries accountability which cannot effectively be applied to large language models \citep{nature2026ai}. The ACL Policy on Publication Ethics explicitly states that generative AI tools may not be listed as authors and that ACL does not consider a generative model capable of fulfilling co-authorship requirements \citep{acl2026policy}.

These policies identify a real issue, but the apparently straightforward concept of ``responsibility'' turns out to cover several different functions. A person may be institutionally responsible for approving a piece of work without possessing all of the detailed expertise embodied in it. Conversely, the person or system best able to explain a particular intellectual decision may have no formal authority over publication.

For the remainder of the paper we therefore avoid treating responsibility as an unanalyzed property. We shift instead to the closely related but more tractable concept of accountability, distinguishing institutional accountability from intellectual accountability. This in turn leads to the more operational notion of \emph{reliable intellectual provenance}: the ability to establish where important ideas and decisions came from, inspect the evidence on which they were based, and obtain an explanation of them.

This leads in turn to a practical question. Can a research project involving substantial AI participation actually be organised so that this kind of provenance is available? Our central technical observation is that the requirements are remarkably similar to those already encountered in serious AI-assisted software development.

But if the technical analogy is as close as we suggest, the original question remains. Why have the two communities developed such different norms? We will argue that part of the explanation may lie not in intrinsic differences between software and research, but in the mechanisms through which their working practices change. Software methodology can evolve directly through individual, team, and organisational practice. Academic research is additionally constrained by formal systems of publication governance whose categories may change on a substantially slower timescale.

Our aim, then, is not to ask the reader to accept ``the vibe'' as an argument. It is to unpack the intuition into a more explicit account of distributed expertise, accountability, provenance, agentic infrastructure, and governance. In that limited sense, we are attempting for \emph{vibe research} something analogous to what the experienced barrister does for Denuto's intuition: to identify more precisely what the practice consists of, why it can work, and how it relates to established institutional concepts.

After developing these arguments, we examine three documented examples from our own work. The present paper is itself being produced through the same methodology, and its development history forms part of the accompanying provenance record.
```
